%%%%%%%%%%%%%%%%%%%%%%%%%%%%%%%%%%%%%%%%%%%%%%%%%%%%%%%%%%%%%%%%%%%%%%%%%%%
% Función: capítulo de ejemplo de contribuciones de la tesis de la tesis por compendio.
%
% Uso: modifica este fichero para adaptar el contenido a tu trabajo.
%%%%%%%%%%%%%%%%%%%%%%%%%%%%%%%%%%%%%%%%%%%%%%%%%%%%%%%%%%%%%%%%%%%%%%%%%%%

\ifthenelse{\equal{\myLanguage}{english}}
{
  \chapter{Thesis Contributions}
}
{
  \chapter{Contribuciones de la tesis}
}
\label{cha:compendium-contributions}

\ifthenelse{\equal{\myLanguage}{english}}
{
  \section{Overview of the Contributions}

  Present and discuss the original contributions of the thesis. Explain how the publications relate to one another, which objective each publication addresses and how their combined results constitute a unified doctoral contribution. The three fictional papers distributed with this example illustrate how to cite each included publication while discussing its contribution~\cite{garcia2026paper1,garcia2026paper2,garcia2026paper3}.

  \subsection{Relationship between the publications}

  Explain the role of each publication and the progression from the initial research question to the combined results of the thesis.

  \section{Supporting Research Material}
  \subsection{Videos and demonstrations}

  Remember that you can very easily include links to videos in the document, and that a list similar to the lists of tables, figures, or algorithms will be generated automatically. The command is defined in \texttt{Config/postamble.tex}, while the list is rendered from \texttt{cover/list-of-videos.tex}. To use it, simply include elements such as \videoLink{Feedback amplifiers: motivation}{https://youtu.be/9xXuasS8B3M} or \videoLink{Feedback amplifiers: ideal theory}{https://youtu.be/ZGMwGX7VCbg}.

  \subsection{Algorithms and source code}

  You can also include algorithms such as Algorithm~\ref{alg:restriction}, and symbols such as~\ac{xidet}.

  In addition, you can include source code listings such as the one shown in Listing~\ref{cod:sample2}.
}
{
  \section{Resumen de las contribuciones}

  Presenta y discute las contribuciones originales de la tesis. Explica cómo se relacionan las publicaciones entre sí, qué objetivo aborda cada una y cómo sus resultados combinados constituyen una contribución doctoral unificada. Los tres artículos ficticios distribuidos con este ejemplo muestran cómo citar cada publicación incluida al discutir su contribución~\cite{garcia2026paper1,garcia2026paper2,garcia2026paper3}.

  \subsection{Relación entre las publicaciones}

  Explica el papel de cada publicación y la progresión desde la pregunta inicial de investigación hasta los resultados conjuntos de la tesis.

  \section{Material de apoyo a la investigación}
  \subsection{Vídeos y demostraciones}

  Recuerda que puedes incluir muy fácilmente enlaces a vídeos en el documento y generar automáticamente un listado similar al de tablas, figuras o algoritmos. La orden se define en \texttt{Config/postamble.tex}, mientras que el listado se genera desde \texttt{cover/list-of-videos.tex}. Para usarla, basta con que incluyas elementos similares a \videoLink{Amplificadores realimentados: motivación}{https://youtu.be/9xXuasS8B3M} o \videoLink{Amplificadores realimentados: teoría ideal}{https://youtu.be/ZGMwGX7VCbg}.

  \subsection{Algoritmos y código fuente}

  Y también puedes incluir algoritmos como el algoritmo~\ref{alg:restriction}, y símbolos como~\ac{xidet}.

  Además de listados de código fuente como en el listado~\ref{cod:sample2}.

}

\begin{algorithm}
  \caption{IntervalRestriction\label{IR}}
  \label{alg:restriction}
  \DontPrintSemicolon
  % \dontprintsemicolon
  \KwData{$G=(X,U)$ such that $G^{tc}$ is an order.}
  \KwResult{$G'=(X,V)$ with $V\subseteq U$ such that $G'^{tc}$ is an
    interval order.}
  \Begin{
    $V \longleftarrow U$\;
    $S \longleftarrow \emptyset$\;
    \For{$x\in X$}{
      $NbSuccInS(x) \longleftarrow 0$\;
      $NbPredInMin(x) \longleftarrow 0$\;
      $NbPredNotInMin(x) \longleftarrow |ImPred(x)|$\;
    }
    \For{$x \in X$}{
      \If{$NbPredInMin(x) = 0$ {\bf and} $NbPredNotInMin(x) = 0$}{
        $AppendToMin(x)$}
    }
    \nl\While{$S \neq \emptyset$}{\label{InRes1}
      \nlset{REM} remove $x$ from the list of $T$ of maximal index\;\label{InResR}
      \lnl{InRes2}\While{$|S \cap ImSucc(x)| \neq |S|$}{
        \For{$ y \in S-ImSucc(x)$}{
          \{ remove from $V$ all the arcs $zy$ : \}\;
          \For{$z \in ImPred(y) \cap Min$}{
            remove the arc $zy$ from $V$\;
            $NbSuccInS(z) \longleftarrow NbSuccInS(z) - 1$\;
            move $z$ in $T$ to the list preceding its present list\;
            \{i.e. If $z \in T[k]$, move $z$ from $T[k]$ to
            $T[k-1]$\}\;
          }
          $NbPredInMin(y) \longleftarrow 0$\;
          $NbPredNotInMin(y) \longleftarrow 0$\;
          $S \longleftarrow S - \{y\}$\;
          $AppendToMin(y)$\;
        }
      }
      $RemoveFromMin(x)$\;
    }
  }
\end{algorithm}


\begin{codefloat}
\lstinputlisting[style=Cnice]{appendix/hello.c}
\caption{Ejemplo de código fuente con estilo \texttt{Cnice}, de nuevo
  con un \texttt{lstinputlisting} dentro de un \texttt{codefloat}}
\label{cod:sample2}
\end{codefloat}

%%% Local Variables:
%%% TeX-master: "../../book"
%%% End:
